\section{Cómo el verifier entra en un ciclo completo de auto mejora}

\subsection{Fase de inferencia}

\begin{itemize}
    \item Best-of-$N$: se genera por completo y luego se hace un rerank;
    \item Beam / tree search: se puntúan los prefijos y se podan;
    \item Early stopping: se detiene la generación cuando la confianza es suficiente;
    \item Critique and revise: se retroalimenta al generator con las razones de la puntuación baja para que revise.
\end{itemize}

\subsection{Fase de entrenamiento}

\begin{itemize}
    \item Filtrado de datos: solo se conservan las synthetic trajectory de puntuación alta;
    \item Rejection sampling fine-tuning: se hace SFT con las trayectorias elegidas por el verifier;
    \item Reinforcement learning: la puntuación del verifier se usa como reward;
    \item Curriculum: los ejemplos de entrenamiento se organizan según la dificultad del paso y el tipo de error.
\end{itemize}

\subsection{Verificación de la estabilidad del ciclo cerrado}

\begin{enumerate}
    \item ¿El generator y el verifier usan conjuntos de evaluación independientes?
    \item Después de actualizar el verifier, ¿los reward antiguos siguen siendo comparables?
    \item ¿El RL eleva la puntuación del verifier pero reduce la tasa de acierto real?
    \item ¿El nuevo generator produce una distribución adversarial que el verifier nunca ha visto?
    \item ¿Se conserva la auditoría humana y un ground truth externo como ancla?
\end{enumerate}

\begin{importantbox}{La coevolución Generator--Verifier debe tener un ancla externa}
Si el generator solo optimiza para el verifier, y el verifier a su vez solo aprende de los datos del generator, ambos pueden ir derivando juntos hacia un protocolo autoconsistente pero erróneo. Las pruebas ejecutables, las reglas formalizadas, la revisión manual por muestreo y los benchmark reservados son la clave para evitar que el ciclo cerrado se distorsione.
\end{importantbox}

\subsection{Resumen de la sección}

El verifier es a la vez un selector, una heurística de búsqueda, un filtro de datos y un reward model. Sus sesgos se amplifican a lo largo del ciclo cerrado, por lo que la tarea de ingeniería más importante en un sistema de auto mejora no es «correr más rondas», sino demostrar que la señal de retroalimentación sigue alineada con el objetivo real.

